% Standalone lecture script. Compile with: latexmk -pdf DQN.tex
\documentclass[11pt,a4paper]{article}
\usepackage[T1]{fontenc}
\usepackage[utf8]{inputenc}
\usepackage{newtxtext,newtxmath}
\usepackage[margin=23mm,top=22mm,bottom=23mm,headheight=15pt]{geometry}
\usepackage{microtype}
\usepackage{amsmath}
\usepackage{xcolor,booktabs,tabularx,enumitem}
\usepackage[most]{tcolorbox}
\usepackage{tikz,pgfplots}
\usepackage{fancyhdr,titlesec}
\usepackage{hyperref}
\pgfplotsset{compat=1.18}
\definecolor{ink}{HTML}{18374A}
\definecolor{accent}{HTML}{087F8C}
\definecolor{pale}{HTML}{EEF6F6}
\definecolor{muted}{HTML}{536572}
\definecolor{warm}{HTML}{BD6735}
\hypersetup{colorlinks=true,linkcolor=accent,urlcolor=accent,citecolor=accent,
  pdftitle={DQN},pdfauthor={RL II 2026}}
\pagestyle{fancy}
\fancyhf{}
\fancyhead[L]{\small\sffamily\color{muted}RL II 2026}
\fancyhead[R]{\small\sffamily\color{muted}DQN}
\fancyfoot[R]{\small\sffamily\color{muted}\thepage}
\renewcommand{\headrulewidth}{0.3pt}
\titleformat{\section}{\Large\sffamily\bfseries\color{ink}}{\thesection}{0.65em}{}
\titleformat{\subsection}{\normalsize\sffamily\bfseries\color{ink}}{}{0pt}{}
\titlespacing*{\section}{0pt}{0pt}{12pt}
\titlespacing*{\subsection}{0pt}{10pt}{4pt}
\titleformat{\subsubsection}{\small\sffamily\bfseries\color{ink}}{}{0pt}{}
\titlespacing*{\subsubsection}{0pt}{10pt}{4pt}
\setlength{\parindent}{0pt}
\setlength{\parskip}{6pt plus 1pt}
\setlist{leftmargin=1.5em,itemsep=4pt,topsep=4pt}
\newtcolorbox{idea}[1]{enhanced,breakable,colback=pale,colframe=accent,
  boxrule=0pt,borderline west={2pt}{0pt}{accent},arc=0pt,
  left=9pt,right=9pt,top=7pt,bottom=7pt,
  title={#1},coltitle=ink,fonttitle=\sffamily\bfseries,
  attach title to upper={\par\smallskip},before skip=9pt,after skip=9pt}
\newcommand{\E}{\mathbb{E}}
\newcommand{\old}{\mathrm{old}}
\newcommand{\clip}{\operatorname{clip}}
\newcommand{\KL}{\operatorname{KL}}
\newcommand{\sg}{\operatorname{sg}}
\newcommand{\Ahat}{\widehat A}
\newcommand{\Vhat}{\widehat V}
\newcommand{\Lclip}{L^{\mathrm{clip}}}
\newcommand{\Lsur}{L^{\mathrm{sur}}}
\newcommand{\nexttopic}{\clearpage}
\begin{document}

{\sffamily\color{accent}\small LECTURE SCRIPT\par}
{\sffamily\bfseries\color{ink}\fontsize{27}{31}\selectfont DQN\par}
\vspace{9pt}
% Based on DQN -- Lecture Notes, Tilman Aach, December 2025.
% At the user's explicit request, original claims and formulas are preserved,
% including the variance identities, parameter-update display, and technical
% claims in the practical sections. New material is marked in source comments.

\section{Training Neural Networks to imitate the Q-function}
% New introductory material requested for this lecture script.
A neural network represents a function $f_\theta$ whose weights and biases are
collected in a parameter vector $\theta$. We cannot generally just set the parameters so that the network represents a
specified function: its outputs depend jointly and nonlinearly on many shared
parameters, and the chosen architecture may only approximate that function.
Instead, we train on
examples $(x_i,y_i)$ with $y_i=f(x_i)$, or on noisy observations of these values.
A \emph{loss function} measures how far a prediction $f_\theta(x_i)$ is from its
target $y_i$. A common choice is the quadratic loss, possibly averaged over a batch $B$:
\begin{equation}
 L_B(\theta)=\frac{1}{|B|}\sum_{i\in B}\bigl(f_\theta(x_i)-y_i\bigr)^2.
\end{equation}
Squaring treats positive and negative errors alike, penalises large errors more
strongly, and gives a smooth loss with a simple derivative. For noisy targets,
minimising expected squared loss recovers the conditional mean. Quadratic loss
is useful, but its sensitivity to large errors can also be a disadvantage.
Backpropagation then computes the gradient of the
loss and an optimiser uses it to adjust the parameters. For example, stochastic
gradient descent (SGD) uses $\theta\leftarrow\theta-\eta\nabla_\theta L_B(\theta)$
with learning rate $\eta>0$. Modern optimisers such as Adam adapt these steps using
gradient history. Each step changes the function represented by the network into the 
direction of the minimum, where the network is as close to the target as possible.

\subsection{Q-learning with a Neural Network}
A Q-network, approximating the Q-function may now take a state as input and outputs a vector of Q-values,
one entry per possible action, i.e., $Q(s,\cdot\,;\theta)$ produces the set
$\{Q(s,a;\theta)\}_{a\in\mathcal A}$. Thinking about how to train this network, we can take two perspectives, for now ignoring the fact
that the functions involved are neural networks and work in the tabular case:
\begin{enumerate}
 \item \textbf{Machine-learning view:} we want to fit a target to the data.
 The Q-function has the property that it is the fixed point of the Bellman optimality operator,
 so at time $n$ we want our function $Q_n(s,a)$ to fit the target $y_n=r+\gamma\max_{a'}Q_n(s',a')$
 as closely as possible for any transitions $(s,a,r,s')$ we observe:
 % Original variance identity retained exactly as requested.
 \begin{equation}\label{eq:original-variance}
 \begin{aligned}
 \widetilde L_n^Q
 &=\E[(\E[y_n]-Q_n(s,a))^2]\\
 &=\E[(y_n-Q_n(s,a))^2]+\E[\mathrm V(y_n)]
 =L_n^Q+\E[\mathrm V(y_n)].
 \end{aligned}
 \end{equation}
 At this point it may not be clear why we want to only minimize $L_n^Q$ in the
 following. However, in the context of DQN it will make sense later.
 \item \textbf{Statistical/RL view:} we want to approximate the Bellman
 optimality operator
 because it is a contraction in the sup-norm, so repeated application converges
 to the optimal $Q$. A well-known statistical trick reveals that
 $\E[X]=\arg\min_\theta\E[(X-\theta)^2]$. Thus, we have
 \begin{equation}
 T^*Q(s,a)=\E_{r,s'}\!\left[r+\gamma\max_{a'}Q(s',a')\right]
 =\arg\min_\theta\E[(y_n-\theta)^2].
 \end{equation}
\end{enumerate}
Doing one-step SGD with stepsize $\frac{\alpha}{2}$ on $L_n^Q$ or on the
minimization problem from the second view gives exactly the same thing,
namely, the Q-Learning update:
\begin{equation}
 Q_{n+1}(s,a)\leftarrow Q_n(s,a)+\alpha\bigl(y-Q_n(s,a)\bigr).
\end{equation}
In the tabular case we can thus extend Q-Learning to a function approximation
problem. The tricks we will see in the next section that characterize DQN and are necessary
to make it work, are usually motivated by the first view, but we think the second 
view is much more natural, specifically for one of the tricks.

\nexttopic
\begingroup
\setlength{\parskip}{5pt plus 1pt}
\setlist{itemsep=2pt,topsep=2pt}
\section{Deep Q-Network (DQN)}
\subsection{Introduction -- Deep Q-Network (DQN)}
Deep Q-Network (DQN), as proposed by Mnih et al. (DeepMind 2015), is
a Q-learning algorithm that uses a deep neural network to approximate the
action-value function. They use the ALE (Arcade learning environment), then
consisting of 49 games (now: 57 games) to show that DQN can match human
level performance.

One of the challenges that DQN addresses is how to deal with a large,
continuous state--action space like that of Atari games: high-dimensional pixel
observations must be processed before Q-learning can be applied. The first
component are convolutional neural networks (CNNs), which process these
images before fully connected layers produce one Q-value per possible action.

\begin{idea}{Problems of Q-learning with Neural Networks}
Before DQN, there were two major probelems with naive function approximation
by NNs that led to instability and divergence in Q-learning:
\begin{itemize}
 \item NNs can forget
 \item Even small Q-updates can massively influence the data distribution
 \item Correlation between successive samples and between the action chosen
 and the target values.
\end{itemize}
\end{idea}

\subsection{DQN Architecture}
The architecture used in the original DQN (``Nature CNN'') is as follows:
\begin{itemize}
 \item \textbf{Pre-processing:} the environment is rendered at $84\times84$
 pixels and converted to grayscale. Four consecutive frames are stacked to
 capture motion. These pixel values together with the score are the input to
 the neural network.
 \item \textbf{Convolutional layers:} three convolutional layers are used.
 The first layer has 32 filters of size $8\times8$ with stride 4, the second
 has 64 filters of size $4\times4$ with stride 2 and the third has 64 filters
 of size $3\times3$ with stride 1. Each convolution is followed by a ReLU
 activation.
 \item \textbf{Fully connected layer:} the flattened convolutional output
 is passed through a fully connected layer with 512 units and ReLU activation.
 \item \textbf{Output layer:} one linear output for each admissible action
 (e.g., 4--18 actions on Atari). Thus the network outputs the vector
 $Q(s,\cdot;\theta)$ for the given state.
\end{itemize}
\begin{center}
{\small
\begin{tabularx}{\linewidth}{@{}lXc@{}}
\toprule
\textbf{Layer} & \textbf{Description} & \textbf{Dimensions}\\
\midrule
Input & Grayscale frames (stack of 4) & $4\times84\times84$\\
Conv1 & $8\times8$ kernel, stride 4, ReLU & $32\times20\times20$\\
Conv2 & $4\times4$ kernel, stride 2, ReLU & $64\times9\times9$\\
Conv3 & $3\times3$ kernel, stride 1, ReLU & $64\times7\times7$\\
Fully Connected & Flatten and dense layer, ReLU & $512$\\
Output & Linear layer for actions & $|\mathcal A|$\\
\bottomrule
\end{tabularx}\par}
\smallskip
{\small\color{muted}Table 1: DQN architecture showing layer descriptions and dimensions.}
\end{center}
\endgroup

\clearpage
\subsection{Experience Replay}
Experience replay addresses the correlation problem between the generated data
and the policy and the fact that training neural networks due to them
forgetting requires batch sampling, which is not possible in the original
Q-Learning algorithm by maintaining a replay buffer $D$ of past transitions.
The algorithm fills the buffer gradually while the agent acts (e.g. using an
$\varepsilon$-greedy policy or any other behavior policy). During training, mini-batches of transitions are
sampled uniformly at random from the buffer. Batch updates are performed
with sampled transitions, and if the buffer is full, the oldest transitions
are discarded. Random sampling and re-use of past experience both reduce
variance and improve data efficiency.

\subsection{Target Networks}
To reduce oscillations and divergence, as well as break the correlation between
the target data distribution and the network to be updated, DQN uses a target
network with parameters $\theta^-$. The target network is updated to the current
network's parameters only every $N$ steps or via a soft update
$\theta^-\leftarrow\tau\theta^-+(1-\tau)\theta$ with $\tau$ close to 1
(although there is no theoretical reason to do so). Between updates,
$\theta^-$ remains fixed, making the target
$y=r+\gamma\max_{a'}Q(s',a';\theta^-)$ independent of the current network's
rapid changes and thereby stabilising training. The above motivation is
grounded in the first view of Q-Learning as fitting your network to a moving
target (that is now made stationary). We can directly see why we only
considered $L_n^Q$ rather than $\widetilde L_n^Q$. The reason is, that now the
variance term is no longer dependend on the parameter to update:
% Original sign and expectations retained exactly as requested.
\begin{equation}
 \widetilde L_n(\theta)
 =\E[(\E[y_n(\theta^-)]-Q_n(s,a;\theta))^2]
 =L_n(\theta)+\E[\mathrm V(y_n(\theta^-))].
\end{equation}
The second view offers a much simpler insight. Fixing the target network for
$N$ steps simply corresponds to applying $N$-step SGD to the minimization problem
\begin{equation}
 \begin{aligned}
 T^*Q(s,a;\theta^-)
 &=\E_{r,s'}\!\left[r+\gamma\max_{a'}Q(s',a';\theta^-)\right]\\
 &=\arg\min_\theta\E[(y_n(\theta^-)-\theta)^2].
 \end{aligned}
\end{equation}
Maintaining a second network doubles memory. Many modern approaches thus
argue to employ other stabilisation techniques, which makes sense if the only
reason to employ target networks is stability. In cases like RL for large modern
models such as LLMs, keeping target networks is prohibitively expensive and
should probably not be done. However, given the second view, we argue that if
one can afford to keep them, one should do so.

There is a slight technicality due to the application of neural networks that
makes the update rule of DQN really differ from tabular Q-Learning. The reason
is that we take the gradient with respect to the weights $\theta$:
% Both original displays (including their factors, signs, left-hand side,
% and unmatched opening parenthesis in the update) are deliberately preserved.
\begin{equation}
 \nabla_\theta l_n(\theta)=\frac12
 (y_n(\theta^-)-Q_n(s,a;\theta))\nabla_\theta Q_n(s,a;\theta)
\end{equation}
leading to
\begin{equation}
 Q_{n+1}(s,a;\theta)=Q_n(s,a;\theta)+\alpha((y_n(\theta^-)-Q_n(s,a;\theta))
 \nabla_\theta Q_n(s,a;\theta).
\end{equation}

\subsection{Algorithm 1: DQN}

{\small
\begin{tabbing}
\hspace{1.3em}\=\hspace{1.3em}\=\kill
\textbf{Initialize} replay buffer $D$, Q-network parameters $\theta$, and target network $\theta^-\leftarrow\theta$;\\
\> \textbf{choose} buffer capacity $|D|$, batch size $|B|$, step size $\alpha$, and target update interval $N$.\\
\textbf{repeat}\\
\> Collect and pre-process a transition $(s,a,s',r)$ using a behavior policy and add it to $D$.\\
\> Sample a mini-batch $B$ of $|B|$ transitions uniformly from $D$.\\
\> Compute targets $y_i=r_i+\gamma\max_{a'}Q_n(s'_i,a';\theta^-)$.\\
\> Take an optimiser step using $\nabla_\theta l_n(\theta)=\frac{1}{2|B|}\sum_{i\in B}(y_i-Q_n(s_i,a_i;\theta))\nabla_\theta Q_n(s_i,a_i;\theta)$.\\
\> Every $N$ steps, update the target network: $\theta^-\leftarrow\theta$.\\
\textbf{until} the training budget is exhausted.
\end{tabbing}}

\nexttopic
\section{Towards Modern Deep RL}
Several algorithmic improvements to DQN, some of which you have already seen in the tabular case, include:
\begin{itemize}
 \item \textbf{Double Q-learning / Double DQN:} Q-learning tends to
 over-estimate action values because the same network is used both to select
 and evaluate the max operator. Double Q-learning maintains two independent
 networks $Q^A$ and $Q^B$. One network selects the greedy action and the other
 evaluates it:
 % The final bias term retains the original coefficient gamma (without alpha).
 \begin{equation}
 \begin{aligned}
 Q^A(s,a)
 &=Q^A(s,a)+\alpha\bigl[r+\gamma Q^B(s',\arg\max_{a'}Q^A(s',a'))-Q^A(s,a)\bigr]\\
 &=Q^A(s,a)+\alpha\bigl[r+\gamma Q^A(s',\arg\max_{a'}Q^A(s',a'))-Q^A(s,a)\bigr]\\
 &\quad+\gamma\bigl(Q^B(s',\arg\max_{a'}Q^A(s',a'))
 -Q^A(s',\arg\max_{a'}Q^A(s',a'))\bigr)
 \end{aligned}
 \end{equation}
 To avoid managing four networks in DQN, Double DQN uses the current
 network to select the best action and the target network to evaluate it,
 thus replacing the full copy by the target network.
 \item \textbf{Clipped double Q / TD3:} Another approach clips the bias-term
 (the latter one in Double Q-Learning) to zero if it is positive, which further
 reduces over-estimation (it explicitly adds the same negative bias but clips
 it to zero if it were to be positive). The deep algorithm (including some other tricks) is called Twin
 Delayed Double Deep Q Networks.
 \item \textbf{Rainbow DQN:} Combining multiple improvements: Double DQN,
 dueling architecture, prioritized replay, noisy nets, multi-step returns and
 distributional RL. There are many further such tricks that can be combined.
\end{itemize}

\subsection{Continuous Action Spaces}
% New discussion prompt requested for this lecture script.
DQN's output layer has one entry for each action. For a continuous
action space, such as a real-valued steering angle or a vector of motor torques,
we cannot allocate an output unit to every possible action. A natural adaptation
is a network that takes both the state and an action as input and outputs a
single value:
\begin{equation}
 (s,a)\longmapsto Q(s,a;\theta)\in\mathbb R.
\end{equation}
\begin{idea}{Why can we not simply do the same as in DQN now?}
Think about what the problem with a DQN-style algorithm would be if we had a continuous action space. Solving this problem will lead us to 
the algorithm Deep Deterministic Policy Gradient (DDPG) and its further developments like Soft Actor-Critic (SAC).
\end{idea}

\subsection{Loss Functions and Normalization}
While the original DQN minimises the squared TD error, the paper mentions
as a remark that clipping the error or using the so-called Huber loss often works
better. The TD error $\delta=r+\gamma\max_{a'}Q(s',a';\theta^-)-Q(s,a;\theta)$
can be clipped into the range $[-1,1]$ before computing the loss: the loss is
quadratic around zero and linear for large errors. The Huber loss is defined as
\begin{equation}
 L_\tau(a)=\begin{cases}
 \frac12 a^2,&\text{if }|a|\leq\tau,\\
 \tau(|a|-\frac\tau2),&\text{otherwise},
 \end{cases}
\end{equation}
where $\tau$ is a threshold hyperparameter and is typically chosen as 1. This
loss becomes the squared loss near zero and the absolute loss for larger errors.
Using this loss helps to prevent very large gradients when the TD error is large
and intuitively helps with stability because large updates are avoided.

Normalizing inputs puts observation components on comparable scales; normalizing
critic outputs and targets consistently reduces sensitivity to reward scale.
Continuous-action actors can output values in $[-1,1]$ and rescale them to the
environment's action bounds.

\nexttopic
\section{Practical Implementation and Evaluation}

To implement DQN using the Stable Baselines3 (SB3) library, key configuration
points include:
\begin{itemize}
 \item \textbf{Feature extractor:} \texttt{policy} selects the architecture.
 SB3's NatureCNN implements the original DQN's convolutional architecture;
 custom feature extractors are also supported.
 \item \textbf{Replay buffer size:} a large buffer (e.g., 1 million
 transitions) is recommended; this requires around 56 GBs of memory.
 \item \textbf{Learning starts:} SB3 delays learning until the buffer has
 been sufficiently filled.
 \item \textbf{Batch size:} number of samples drawn from the buffer per
 gradient step.
 \item \textbf{Soft/hard target update:} SB3 supports periodic hard updates
 or soft updates controlled by $\tau$ ($\tau=1$ gives a hard update and is the default).
 \item \textbf{\texttt{target\_update\_interval}:} number of gradient steps between hard updates.
 \item \textbf{\texttt{train\_freq} and \texttt{n\_envs}:} train every
 \texttt{train\_freq} environment steps, collecting each batch from all
 \texttt{n\_envs} parallel environments.
 \item \textbf{Multi-step returns:} SB3 supports multi-step targets and
 trajectory sampling, which correspond to $n$-step Bellman updates.
 \item \textbf{Huber loss and clipping:} SB3 can apply Huber loss and clip
 the TD error when computing the loss.
\end{itemize}

\subsection{Atari Benchmarks and Evaluation Metrics}
Atari evaluation metrics compared DQN's scores to two baselines: a random
agent and a human expert. The human-normalized score (HNS) was defined as
$(\mathrm{agent}-\mathrm{random})/(\mathrm{human}-\mathrm{random})$.
Comparisons can be unreliable because some games permit extremely large scores
and random seeds can drastically affect results. Robust metrics like the
inter-quartile mean (IQM) or median across many seeds mitigate these issues.
The rliable package provides these metrics and bootstrap confidence intervals.

Benchmarks include the Atari 200M (million frames) or 1B (billion frames),
where the number of steps is the same number but divided by 4 (4 frames are
stacked). There are also papers for example giving the 1/3/5/10 most predictive games in
order to get a more valid analysis without testing all games.

To make Atari environments less deterministic and to encourage robust policies,
several forms of environment randomisation are used:
\begin{itemize}
 \item \textbf{Noisy networks (NoisyNets):} learned noise in network weights
 makes Q-values stochastic, replacing $\varepsilon$-greedy exploration with
 parameter noise without randomising the environment itself.
 \item \textbf{Sticky actions:} probabilistically repeat the previous action
 despite a new request. Proposed in 2017 after DQN, this helps prevent
 exploitation of deterministic dynamics.
 \item \textbf{Random frame skipping:} instead of always repeating each
 action for a fixed number of frames (e.g., 4), choose a number of skipped
 frames randomly in a range.
 \item \textbf{Initial no-ops:} start each episode with a random number of
 no-op actions to randomise the initial state. A fixed number still gives a deterministic game.
 \item \textbf{Random action noise:} occasionally execute a random action
 instead of the policy's action. This is mostly discouraged, as it leads to high
 instability.
\end{itemize}
These techniques make the training environment more stochastic and reduce
the risk of over-fitting to deterministic trajectories.

\end{document}
